\section{FR Representations and Transformations}
\label{sec:operators}

This section does not study screening. It constructs the mathematical
universe in which screening will later be defined: Fenchel--Rockafellar
representations together with their primitive transformations. To
study screening structurally, we first introduce a class of
optimization representations together with the primitive
transformations acting on them. These representations are closed under
the transformations considered in this paper, allowing screening to be
formulated as a transformation of representations rather than as a
model-specific optimization procedure.

The section answers exactly one question -- what is the mathematical
universe in which screening will be studied -- in three tiers: the
object (an FR representation), the transformations acting on it including duality, translation, restriction, and the structural law each transformation
satisfies. 

%%%%%%%%%%%%%%%%%%%%%%%%%%%%%%%%%%%%%%%%%%%%%%%%%%%%%%%%%%%%%%%

\subsection{FR Representations}

\paragraph{Motivation.}
Every object in this paper -- a convex optimization problem, its
translate, its restriction, its dual -- is encoded as a single
algebraic quadruple, so that "translate," "restrict," and "dualize"
can each be defined once, uniformly, instead of separately for every
problem class that happens to arise.


We consider the problem of the form
\[
\min_{x \in \mathbb R^n}\;
p(x), \qquad p(x) = f(Ax)+g(x)+a.
\]
where
\(
f:\mathbb R^m\rightarrow\overline{\mathbb R},
g:\mathbb R^n\rightarrow\overline{\mathbb R},
A\in\mathbb R^{m\times n},
a\in\mathbb R,
\)
with \(f\) and \(g\) closed, proper, and convex.
We further assume that \(f\) and \(g\) are separable, \(f(y)=\sum_{j=1}^m f_j(y_j)\)
and \(g(x)=\sum_{i=1}^n g_i(x_i)\), with each \(f_j\) and \(g_i\) closed, proper, and convex on $\mathbb{R}$. Here, at the first glance, \(a\) is a redundant constant since it does not affect the minimizer of \(p\). However, it is important to keep track of \(a\) because it is used as an absorbing constant in the following transformations. 

The pair \((m,n)\) is called the \emph{problem size}, also the \emph{intrinsic complexity} of the problem. 
We define $\Gamma_{m, n}$ to be the set of all function $p$ with problem size $(m, n)$, and $\Gamma$ to be the set of all FR representations of any size.

\begin{definition}[FR representation]
A Fenchel--Rockafellar representation is a quadruple
\[
p\equiv (f,g,A,a),
\]
\end{definition}




% \noteLE{Regularity added 2026-07-03 (paper.md TODO): every identity
% below that involves a conjugate needs closed, proper, convex functions
% for biconjugation (\(h^{**}=h\)) to hold. The Duality theorem in
% Section~\ref{sec:screening} needs this not just for \(f\), \(g\)
% themselves but for the separable blocks \(f_{\bar J}\), \(g_{\bar I}\)
% that result from restriction -- see Lemma~\ref{lem:separable-sum-regular}
% below for why those inherit it too.}



An FR representation is not merely a convenient way to write down a
problem. The quadruple \(p=(f,g,A,a)\) is designed so that the class of
FR representations is \emph{closed} under the three transformations
studied in this section including duality, translation and restriction.

%%%%%%%%%%%%%%%%%%%%%%%%%%%%%%%%%%%%%%%%%%%%%%%%%%%%%%%%%%%%%%%

\subsection{Duality}

% main message: the dual of an FR representation is another FR representation, and duality is an involution on the class of FR representations. 


The Fenchel--Rockafellar dual of \(p=(f,g,A,a)\) is
\[
\max_{y \in \mathbb R^m}\;
-d(y), \qquad d(y) = g^*(-A^\top y) + f^*(y)-a.
\]
Here the constant \(a\) is negated in the dual, so that the weak duality inequality \(p(x) + d(y)\geq 0\) remains valid for all \(x\in \mathbb R^n\) and \(y \in \mathbb R^m\).
The dual function $d$ is itself an FR representation, with the problem size \((n, m)\). Therefore, FR duality is a transformation on the class of FR representations. 


\begin{definition}[Duality]
The Fenchel--Rockafellar dual of \(p=(f,g,A,a)\) is
\[
p^*
:=
(g^*,\,f^*,\,-A^\top,\,-a).
\]
\end{definition}

Now, $*:\Gamma \to \Gamma$ is a transformation on the class of FR representations. 

Duality exchanges the roles of \(f\) and \(g\), and it exchanges domain
translation with slope translation (a function's domain and its slope
are conjugate to each other, in the same sense that \(f\) and \(g^*\)
are). The value parameter \(c\) plays no such role: it shifts the
representation's own constant \(a\) directly, and duality negates that
constant (just defined above). So \(k=(b,c,d)\) should dualize to a
translation parameter that swaps and negates \(b\) and \(d\), and also
negates \(c\), to stay consistent with the constant's own sign flip.



\begin{proposition}[Involution]
\label{prop:involution}
For every FR representation \(p \in \Gamma\), we have \((p^*)^*=p\).
\end{proposition}

\begin{proof}
Write \(p=(f,g,A,a)\), so \(p^*=(g^*,f^*,-A^\top,-a)\). Applying the
same rule again,
\[
(p^*)^*
=
\bigl((f^*)^*,\,(g^*)^*,\,-(-A^\top)^\top,\,-(-a)\bigr)
=
(f^{**},g^{**},A,a)
=
(f,g,A,a)
=
p,
\]
using biconjugation \(f^{**}=f\), \(g^{**}=g\), which hold true since $f$ and $g$ are closed, proper, and convex.
\end{proof}

\noteLE{This was missing before 2026-07-03; the Simultaneous Screening
corollary of Section~\ref{sec:screening} uses it directly in its proof.
Note it holds regardless of whether the dual's constant is \(a\) or
\(-a\) -- negating twice returns to the original either way -- so this
proposition alone could not have told us which convention to pick; the
Duality theorem is what forced \(-a\) (see the correction below).}









%%%%%%%%%%%%%%%%%%%%%%%%%%%%%%%%%%%%%%%%%%%%%%%%%%%%%%%%%%%%%%%


\subsection{Translation}

\paragraph{Motivation.}
A screening operator will eventually need to absorb the contribution of
a coordinate it is about to drop into the coordinates it keeps. That
absorption happens at the level of individual convex functions, before
it is ever assembled into an operation on a full representation -- hence
three primitive translations, one for each way a function can be
shifted: in its domain, in its slope, in its value. All three are the
same kind of object -- a translation -- so they are stated as one
family, distinguished only by which part of the function they shift.

\begin{definition}[Primitive translations]
For \(b\in\mathbb R^n\), \(c\in\mathbb R\), and \(d\in\mathbb R^n\),
define, for a function \(h: \mathbb R^n \rightarrow \overline{\mathbb R}\),
\[
\text{domain translation:}
\quad
(T^{\mathrm d}_bh)(y):=h(y+b),
\]
\[
\text{value translation:}
\quad
(T^{\mathrm v}_ch)(y):=h(y)+c,
\]
\[
\text{slope translation:}
\quad
(T^{\mathrm s}_dh)(x):=h(x)+\langle d,x\rangle.
\]
\end{definition}

\paragraph{Insight.}
Each of the three is a group action of its parameter space on the
functions of the relevant type: composing two domain translations adds
their parameters, and likewise for value and slope translations.

% \begin{proposition}[Abelian primitive translations]
% Each translation family
% \[
% \{T^{\mathrm d}_b\},
% \qquad
% \{T^{\mathrm v}_c\},
% \qquad
% \{T^{\mathrm s}_d\}
% \]
% forms an Abelian family under composition.
% \end{proposition}

% \noteLE{Not currently cited by a Main Theorem. It was written for the
% now-deleted Commutativity theorem of Section~\ref{sec:screening}; kept
% as a basic structural fact about the primitive translations.}

These three per-function translations combine into a single operation
on a full FR representation.

\begin{definition}[Translation parameter and its dual]
    Let $K= \mathbb{R}^m \times \mathbb{R} \times \mathbb{R}^n$ be the space of translation parameters. A translation parameter is \(k=(b,c,d)\in K\). 
    Its dual parameter is defined as
    \[
    k^*
    :=
    (-d,\;-c,\;-b).
    \]
\end{definition}

Note that $K$ is an Abelian group under addition, and the duality map $k \mapsto k^*$ is an involution on $K$.

% \noteLE{\textbf{History, updated 2026-07-03.} This definition has been
% corrected twice, in opposite directions, and it is worth recording why.
% Originally \(k^*=(-d,-c,-b)\), matched with \(p^*=(g^*,f^*,-A^\top,a)\)
% (constant unchanged) -- inconsistent, since \((T_k(p))^*\) then has
% constant \(a+c\) while \(T_{k^*}(p^*)\) has constant \(a-c\). First fix
% (same day, earlier pass): keep \(p^*\) with \(a\) unchanged, change
% \(k^*\) to \((-d,c,-b)\). This worked for Fenchel Equivariance alone,
% but forced sample screening's canonical constant to \(+f_{\bar J}^*(s)\)
% in Section~\ref{sec:screening}, which conflicted with the independent
% optimization-motivated derivation (Fenchel--Young equality), which gives
% \(-f_{\bar J}^*(s)\). Second fix (this version): negate \(p^*\)'s
% constant instead (\(p^*=(g^*,f^*,-A^\top,-a)\)), which restores the
% \emph{original} \(k^*=(-d,-c,-b)\) and, checked independently, also makes
% the two sample-screening derivations agree. This is the version used
% throughout Section~\ref{sec:screening} now.}



\begin{definition}[Translation action]
Let \(p=(f,g,A,a)\). The action of \(k=(b,c,d)\) on \(p\) is defined by
\[
T_k(p)
:=
(T^{\mathrm d}_bf,\;
T^{\mathrm s}_dg,\;
A,\;
a+c).
\]
\end{definition}

\paragraph{Insight.}
The matrix \(A\) is untouched by translation: only the two functions
and the constant move. This is what lets restriction and translation
later be composed without the matrix itself being disturbed by which
primitive acts first.

\begin{proposition}[Abelian translation action]
The family \(\{T_k\}\) forms an Abelian translation action on
the class of FR representations:
\[
T_{k_1+k_2}(p)
=
T_{k_1}(T_{k_2}(p)).
\]
\end{proposition}

\noteLE{Not used by any result in Section~\ref{sec:screening}. Kept as
a basic structural fact: it is what makes \(T_k\) an honest group
action.}

Section~\ref{sec:screening} will apply the translation action after
choosing \(k\) as a function of the coordinates being screened. Before
that, we establish this subsection's law: how the translation action
interacts with Fenchel--Rockafellar duality.



\begin{proposition}
\label{prop:domain-dual}
For every convex function $f: \mathbb{R}^n \to \overline{\mathbb{R}}$, $b,d \in \mathbb{R}^n$, and $c \in \mathbb{R}$, we have 
\begin{enumerate}
    \item \((T^{\mathrm d}_bf)^* = T^{\mathrm s}_{-b}(f^*)\)
    \item \((T^{\mathrm s}_dg)^* = T^{\mathrm d}_{-d}(g^*)\)
    \item \((T^{\mathrm v}_cf)^* = T^{\mathrm v}_{-c}(f^*)\)
\end{enumerate}
\end{proposition}

\noteLE{Used in the Duality theorem (Section~\ref{sec:screening}) and
in the proof of Fenchel Equivariance below.}

\noteLE{Not currently cited. This identity is about conjugating a
value-\emph{shifted function}; it is not what governs how the
translation action's own \(c\)-parameter behaves under duality -- that
is elementary (see the proof below) and does not go through
\(T^{\mathrm v}_c\)/\(T^{\mathrm v}_{-c}\) at all. Kept for completeness,
as the third of the three primitive identities.}

The previous identities admit a unified formulation -- this
subsection's law.

\paragraph{Statement (the law of translation).}

\begin{theorem}[Fenchel Equivariance]
\label{thm:fenchel-equivariance}
Let \(p\) be an FR representation and let \(k=(b,c,d)\) be a translation
parameter. Then
\[
\boxed{
(T_k(p))^*
=
T_{k^*}(p^*).
}
\]
\end{theorem}

\paragraph{Proof.}
Write \(p=(f,g,A,a)\), so \(T_k(p)=(T^{\mathrm d}_bf,\,T^{\mathrm
s}_dg,\,A,\,a+c)\). Applying the FR dual, which negates the constant
slot, \((f_0,g_0,A_0,a_0)^*=(g_0^*,f_0^*,-A_0^\top,-a_0)\), termwise, and
using Proposition~\ref{prop:domain-dual},
\[
(T_k(p))^*
=
\big(
(T^{\mathrm s}_dg)^*,\;
(T^{\mathrm d}_bf)^*,\;
-A^\top,\;
-(a+c)
\big)
=
\big(
T^{\mathrm d}_{-d}(g^*),\;
T^{\mathrm s}_{-b}(f^*),\;
-A^\top,\;
-a-c
\big).
\]

On the other side, \(p^*=(g^*,f^*,-A^\top,-a)\), and \(k^*=(-d,-c,-b)\),
so \(T_{k^*}(p^*)\) applies a domain translation by \(-d\) to \(p^*\)'s
first slot, a slope translation by \(-b\) to its second slot, and adds
\(-c\) to its constant:
\[
T_{k^*}(p^*)
=
\big(
T^{\mathrm d}_{-d}(g^*),\;
T^{\mathrm s}_{-b}(f^*),\;
-A^\top,\;
-a-c
\big).
\]
Every term matches \((T_k(p))^*\), which proves the theorem.
\qed

\paragraph{Interpretation.}
Translation is not just a transformation on FR representations -- it is
compatible with Fenchel--Rockafellar duality. Duality acts on the whole
translation \emph{action}, not just on the representation: it is the
same to translate first and dualize, or to dualize first and translate
by the dual parameter. This is why Section~\ref{sec:screening} can
define feature screening and sample screening independently and still
get an exact duality between them: both are built from this one
equivariant action, applied to the two different restrictions of the
next subsection.

\noteLE{Updated 2026-07-05: now cited by name. The Dual of screening
theorem in Section~\ref{sec:screening} was rewritten to combine
Decomposition, this theorem, and Restriction Equivariance
(Proposition~\ref{prop:restriction-equivariance}) as three explicit
steps, instead of re-deriving the equivalent computation directly from
Proposition~\ref{prop:domain-dual} and the restriction lemmas alone.
Those two lemmas are still used, but only to verify that the dual
translation parameter produced by this theorem coincides with the one
that directly defines screening on \(p^*\).}

%%%%%%%%%%%%%%%%%%%%%%%%%%%%%%%%%%%%%%%%%%%%%%%%%%%%%%%%%%%%%%%









%%%%%%%%%%%%%%%%%%%%%%%%%%%%%%%%%%%%%%%%%%%%%%%%%%%%%%%%%%%%%%%








\subsection{Restriction}

\paragraph{Motivation.}
Translation alone only ever moves a coordinate's contribution around;
it never removes it. The second primitive, coordinate restriction, is
what actually drops a block of coordinates from the representation --
this is the other half of Screening~=~Translation~+~Restriction.


Let \(I\subseteq\{1,\ldots,n\}\) and \(J\subseteq\{1,\ldots,m\}\). For a
separable function \(g(x)=\sum_{i=1}^n g_i(x_i)\), write
\(x_I:=(x_i)_{i\in I}\) and \(g_I(x_I):=\sum_{i\in I}g_i(x_i)\); define
\(f_J\) symmetrically for a separable \(f(y)=\sum_{j=1}^m f_j(y_j)\).
For a matrix \(A\in\mathbb R^{m\times n}\), write \(A_I:=A_{:,I}\) for
the submatrix of columns indexed by \(I\), and \(A_J:=A_{J,:}\) for the
submatrix of rows indexed by \(J\). Combining both,
\(A_{J,I}:=(A_J)_I=(A_I)_J\) denotes the submatrix with rows in \(J\)
and columns in \(I\); this block notation is what makes the
Simultaneous Screening discussion in Section~\ref{sec:screening}
precise about which cross term appears.

\begin{definition}[Feature and sample restriction]
Let \(I\sqcup\bar I=\{1,\ldots,n\}\). Define
\[
R^{F}_I(f,g,A,a)
:=
(f,g_I,A_I,a).
\]
Let \(J\sqcup\bar J=\{1,\ldots,m\}\). Define
\[
R^{S}_J(f,g,A,a)
:=
(f_J,g,A_J,a).
\]
\end{definition}

\paragraph{Insight.}
Feature restriction only touches \(g\) and the columns of \(A\); sample
restriction only touches \(f\) and the rows of \(A\). Because they act
on disjoint parts of \(p\), the order in which they are applied cannot
matter -- this subsection's law.


\begin{proposition}[Restriction Equivariance]
\label{prop:restriction-equivariance}
For every \(p=(f,g,A,a)\in\Gamma\), every \(I\subseteq\{1,\ldots,n\}\),
and every \(J\subseteq\{1,\ldots,m\}\),
\[
\bigl(R^{F}_I(p)\bigr)^*=R^{S}_I(p^*),
\qquad
\bigl(R^{S}_J(p)\bigr)^*=R^{F}_J(p^*).
\]
\end{proposition}

\begin{proof}
Write \(p^*=(g^*,f^*,-A^\top,-a)\). Applying the FR dual termwise to
\(R^F_I(p)=(f,g_I,A_I,a)\),
\[
\bigl(R^F_I(p)\bigr)^*
=
\bigl((g_I)^*,\;f^*,\;-A_I^\top,\;-a\bigr).
\]
By separability of \(g\), \((g_I)^*=(g^*)_I\). Rows \(I\) of
\(-A^\top\) equal \(-A_I^\top\), since rows of \(A^\top\) are columns of
\(A\) transposed. Hence
\[
\bigl(R^F_I(p)\bigr)^*
=
\bigl((g^*)_I,\;f^*,\;-A_I^\top,\;-a\bigr)
=
R^S_I(p^*),
\]
the last equality by the definition of sample restriction applied to
\(p^*=(g^*,f^*,-A^\top,-a)\). The second identity is the mirror
computation, restricting \(R^S_J(p)=(f_J,g,A_J,a)\) instead, and using
\((f_J)^*=(f^*)_J\) from the same lemma.
\end{proof}

\noteLE{Added 2026-07-05: restriction's counterpart to Fenchel
Equivariance. Translation's law says dualizing a translation gives the
same translation, dual parameter, on the dual representation
(\(k\mapsto k^*\)); this says dualizing a feature restriction gives
exactly a sample restriction, \emph{same index set}, on the dual
representation, and symmetrically. Updated same day: now cited by name
in Section~\ref{sec:screening}'s Duality theorem, rewritten to combine
Decomposition, Fenchel Equivariance, and this proposition as three
explicit steps rather than one monolithic computation; the
restriction-conjugate lemmas are still used, but only to verify the
resulting translation parameter. Positioning relative to Commutativity
(this subsection's other law) is still open: paper.md's TODO on
Duality as a third transformation should say whether this or
Commutativity is "the" law of this subsection, or whether both are,
now that duality is treated as a transformation on the same footing as
translation and restriction.}

\begin{proposition}[Commutativity]
\label{prop:restrictions-commute}
Feature restriction and sample restriction commute:
\[
\boxed{
R^{F}_I
R^{S}_J
=
R^{S}_J
R^{F}_I.
}
\]
\end{proposition}

\noteLE{Restriction's structural counterpart to Fenchel Equivariance --
repositioned 2026-07-05 from an incidental, unused proposition to this
subsection's law, on the same footing as translation's law in the
previous subsection. Still stated without proof (kept without proof
per author instruction, 2026-07-03); the proof is one line, immediate
from the Insight above (disjoint parts of \(p\)), and adding it is a
separate, still-open decision (paper.md TODO), not resolved by this
repositioning. Not currently cited by any result of
Section~\ref{sec:screening} -- would be needed if a future result
composes \(S^F\) and \(S^S\) on the same representation and compares
orders directly (the dropped Commutativity theorem tried exactly this
and failed for a different reason, an interaction through \(A\), not
because this proposition is false; see paper.md's Open Problems
history).}

\begin{proposition}[Restriction is compatible with translation]
\[
\boxed{
R(T_k(p))
=
T_{k'}(R(p)),
}
\]
where \(k'\) is obtained by restricting the corresponding translation
parameters to the remaining coordinates.
\end{proposition}

\noteLE{Kept without proof, per author instruction (2026-07-03). Not
currently cited by a result in Section~\ref{sec:screening}. It was
written for the now-deleted Commutativity theorem, which turned out to
need more than this -- the gap found there was not in this proposition
itself (a single restriction after a single translation, on one
representation, does commute in this sense) but in composing \emph{two
different} screening operators whose translation parameters are each
derived from \(p\)'s own data. Kept as a true, general fact.}

\begin{lemma}[Conjugate of a separable restriction]
\label{lem:restrict-conjugate}
Let \(h:\mathbb R^n\rightarrow\overline{\mathbb R}\) be separable,
\(h(x)=\sum_{i=1}^n h_i(x_i)\), and let \(I\sqcup\bar I=\{1,\ldots,n\}\).
Writing \(h_I(x_I):=\sum_{i\in I}h_i(x_i)\),
\[
\boxed{
(h_I)^*=(h^*)_I.
}
\]
\end{lemma}

\begin{proof}
Because \(h\) is separable, the supremum defining the conjugate splits
over the independent coordinates \(x_i\):
\(h^*(y)=\sum_{i=1}^n h_i^*(y_i)\). Restricting the same supremum to the
block \(I\) alone,
\[
(h_I)^*(y_I)
=
\sup_{x_I}\Big\{\langle y_I,x_I\rangle-\sum_{i\in I}h_i(x_i)\Big\}
=
\sum_{i\in I}h_i^*(y_i)
=
(h^*)_I(y_I).
\]
\end{proof}

\begin{lemma}[Separable sums preserve regularity]
\label{lem:separable-sum-regular}
Let \(h=\sum_{i=1}^n h_i\) be separable with each \(h_i\) closed, proper,
convex. Then for every \(I\subseteq\{1,\ldots,n\}\), the restricted
function \(h_I=\sum_{i\in I}h_i\) is also closed, proper, convex.
\end{lemma}

\begin{proof}
Convexity is immediate: a sum of convex functions is convex. For
properness, each \(h_i\) has a point where it is finite and is nowhere
\(-\infty\); choosing such a point in each coordinate \(i\in I\) gives a
point where \(h_I\) is finite, and \(h_I\) is nowhere \(-\infty\) since
no summand is. For closedness, a finite sum of closed (lower
semicontinuous) functions is closed: if \(x^{(k)}\to x_I\), then
\(\liminf_k h_I(x^{(k)})=\liminf_k\sum_{i\in I}h_i(x^{(k)}_i)\geq
\sum_{i\in I}\liminf_k h_i(x^{(k)}_i)\geq\sum_{i\in I}h_i(x_i)=h_I(x_I)\),
using lower semicontinuity of each \(h_i\) and superadditivity of
\(\liminf\) over a finite sum.
\end{proof}

\noteLE{Added 2026-07-03 (paper.md TODO). Needed by the Duality theorem
of Section~\ref{sec:screening}: that proof requires \(g_{\bar I}\)
(resp. \(f_{\bar J}\)) itself closed, proper, convex, to invoke
biconjugation \(g_{\bar I}^{**}=g_{\bar I}\). The Separable
representation definition above only states this for the individual
summands \(g_i\); this lemma is what lets it transfer to any sub-sum.}

This section introduced the mathematical universe used throughout the
paper: a class of Fenchel--Rockafellar representations together with
two primitive transformations, translation and restriction. The two
transformations satisfy their respective structural laws -- Fenchel
Equivariance and Commutativity -- which provide the mathematical
ingredients used in the next section to define screening and establish
Feature--Sample Duality.
